% MR47091C
% Tachykarde Attacke
% 14.0
% 2013-11-30

\label{m:tachykarde-attacke}
\EE{Taktik}: \Speech{\Ca{Vitale Bedrohung!}
Ärztliche Therapie sinnvoll.}
\begin{itemize}
    \item Lagerung: \EE{Oberkörper hoch}
  \item Striktes \EE{Bewegungsverbot!}
  \item Beengende Kleidung öffnen
    \item \ce{O2}-Gabe gemäß \prettyref{m:sauerstoffberieselung}
  \item \TxBeiVit \TxMassVitMKBes
  \begin{itemize}
    \item Notarzt: Eine medikamentöse Therapie ist meistens
    erforderlich.
  \end{itemize}
  \item Psychische Betreuung! Oase der Ruhe schaffen.
  \item Wenn der Patient versorgt ist (d.\,h. alle obigen Punkte erfüllt
  wurden) soll -- wenn möglich -- ein \E{EKG}
  (Extremitätenableitung, \prettyref{topic:ekg}) abgeleitet werden. Der
  Ausdruck dient v.\,a. der \E{Dokumentation} und der Information für das
  nachbehandelnde Personal (Notarzt, Spital, Facharzt, \ldots).
\end{itemize}